%% По требованиям ИТМО итоговая аннотация формируется в ИСУ и подшивается
%% отдельно. Данный раздел оставлен для рабочей версии текста.

\structuralelement{Аннотация}

\noindent
В работе решается задача разработки программного обеспечения автономного
мобильного омни-робота, обеспечивающего сопровождение пользователя в типовой
квартирной среде при наличии динамических препятствий. Описан полный контур
системы: от воспроизводимой контейнерной симуляционной среды на базе
Docker, ROS~2 и Gazebo до алгоритма управления локального планирования на
основе стохастической реализации модельного прогнозирующего управления
(Model Predictive Path Integral, MPPI). Приведена математическая модель
омни-платформы с колесами типа Mecanum, сформирована функция потерь,
формализующая одновременно требования по удержанию дистанции до цели,
центрированию цели в поле зрения камеры, безопасному обходу динамических
препятствий и корректному поведению при потере цели. 

Практическая часть включает описание архитектуры программного
комплекса, реализации двух пользовательских критиков
MPPI-контроллера~--- центрирования цели в кадре камеры
(\texttt{TargetCameraCritic}) и направленного штрафа поля препятствий
(\texttt{PotentialFieldCritic})~--- а также автоматизированной серии
из 40 безголовых запусков с количественной оценкой режимов
сопровождения, видимости цели, безопасности относительно препятствий
и плавности управления.